\documentclass[11pt,openany]{report}
\usepackage[letterpaper,margin=0.82in,headheight=14pt]{geometry}
\usepackage[T1]{fontenc}
\usepackage[utf8]{inputenc}
\usepackage{lmodern}
\usepackage{microtype}
\usepackage{amsmath,amssymb,amsthm,mathtools}
\usepackage{booktabs,longtable,tabularx,array,multirow}
\usepackage{graphicx,float,placeins}
\usepackage[dvipsnames]{xcolor}
\usepackage{enumitem}
\usepackage{fancyhdr}
\usepackage{listings}
\usepackage{url}
\usepackage[hidelinks]{hyperref}
\usepackage[nameinlink,noabbrev]{cleveref}
\usepackage{needspace}
\allowdisplaybreaks

\hypersetup{
  pdftitle={Projector-Generated Special Functions II},
  pdfauthor={Komorebi P.},
  pdfsubject={Cyclotomic Stieltjes packets, positive energy resolvents, Gamma scaling, and character decoherence}
}

\definecolor{Ink}{HTML}{17202A}
\definecolor{Blue}{HTML}{2E86AB}
\definecolor{Teal}{HTML}{2A9D8F}
\definecolor{Orange}{HTML}{F4A261}
\definecolor{Magenta}{HTML}{A23E8C}
\definecolor{Red}{HTML}{E76F51}
\definecolor{Paper}{HTML}{FAFAF7}
\definecolor{Muted}{HTML}{5D6D7E}
\definecolor{CodeBG}{HTML}{F3F5F5}

\pagecolor{Paper}
\color{Ink}
\graphicspath{{figures/}}
\setlength{\parindent}{0pt}
\setlength{\parskip}{0.52em}
\raggedbottom

\pagestyle{fancy}
\fancyhf{}
\fancyhead[L]{\small Projector-Generated Special Functions II}
\fancyhead[R]{\small Komorebi P.}
\fancyfoot[C]{\thepage}
\renewcommand{\headrulewidth}{0.4pt}

\newtheorem{theorem}{Theorem}[chapter]
\newtheorem{proposition}[theorem]{Proposition}
\newtheorem{lemma}[theorem]{Lemma}
\newtheorem{corollary}[theorem]{Corollary}
\newtheorem{conjecture}[theorem]{Conjecture}
\theoremstyle{definition}
\newtheorem{definition}[theorem]{Definition}
\newtheorem{example}[theorem]{Worked example}
\newtheorem{algorithm}[theorem]{Algorithm}
\theoremstyle{remark}
\newtheorem{remark}[theorem]{Remark}
\newtheorem{boundary}[theorem]{Boundary statement}

\newcommand{\C}{\mathbb C}
\newcommand{\R}{\mathbb R}
\newcommand{\N}{\mathbb N}
\newcommand{\Z}{\mathbb Z}
\newcommand{\Q}{\mathbb Q}
\newcommand{\e}{\mathrm e}
\newcommand{\ii}{\mathrm i}
\newcommand{\Bfrak}{\mathfrak B}
\newcommand{\Wcal}{\mathcal W}
\newcommand{\Vcal}{\mathcal V}
\newcommand{\Ncal}{\mathcal N}
\newcommand{\Ecal}{\mathcal E}
\newcommand{\Mcal}{\mathcal M}
\newcommand{\Ccal}{\mathcal C}
\newcommand{\norm}[1]{\left\lVert#1\right\rVert}
\newcommand{\abs}[1]{\left\lvert#1\right\rvert}
\newcommand{\dd}{\,\mathrm d}
\newcommand{\diag}{\operatorname{diag}}
\newcommand{\Res}{\operatorname{Res}}
\newcommand{\Corr}{\operatorname{Corr}}
\newcommand{\Var}{\operatorname{Var}}
\newcommand{\E}{\mathbb E}
\newcommand{\ket}[1]{\left|#1\right\rangle}
\newcommand{\bra}[1]{\left\langle#1\right|}
\newcommand{\braket}[2]{\left\langle#1\middle|#2\right\rangle}

\newcolumntype{L}[1]{>{\raggedright\arraybackslash}p{#1}}

\lstdefinestyle{pythonstyle}{
  language=Python,
  basicstyle=\ttfamily\scriptsize,
  backgroundcolor=\color{CodeBG},
  frame=single,
  rulecolor=\color{Muted!45},
  breaklines=true,
  breakatwhitespace=false,
  showstringspaces=false,
  numbers=left,
  numberstyle=\tiny\color{Muted},
  numbersep=7pt,
  tabsize=4,
  columns=fullflexible,
  keepspaces=true
}
\lstset{style=pythonstyle}

\newcommand{\FigureWide}[3]{%
\Needspace{0.53\textheight}
\begin{figure}[H]
\centering
\includegraphics[width=0.92\textwidth,height=0.44\textheight,keepaspectratio]{#1.pdf}
\caption{#2 #3}
\label{fig:#1}
\end{figure}
}

\newcommand{\InlinePart}[1]{%
\FloatBarrier
\Needspace{10\baselineskip}
\refstepcounter{part}%
\addcontentsline{toc}{part}{\protect\numberline{\thepart}#1}%
\vspace{1.4em}
{\color{Blue}\Large\bfseries Part \thepart\par}
\vspace{0.2em}
{\color{Blue}\LARGE\bfseries #1\par}
\vspace{0.9em}
}

\newcommand{\FlowChapter}[1]{%
\FloatBarrier
\Needspace{14\baselineskip}
\refstepcounter{chapter}%
\setcounter{section}{0}%
\setcounter{subsection}{0}%
\addcontentsline{toc}{chapter}{\protect\numberline{\thechapter}#1}%
\markboth{#1}{#1}%
\vspace{1.3em}
{\color{Muted}\large\bfseries Chapter \thechapter\par}
\vspace{0.15em}
{\LARGE\bfseries #1\par}
\vspace{0.9em}
}

\title{\textbf{Projector-Generated Special Functions II}\\[0.25em]
\Large Cyclotomic Stieltjes Packets, Packet Resultants, Positive Energy Resolvents, Gamma Scaling, and Character Decoherence}
\author{Komorebi P.}
\date{}

\begin{document}
\maketitle

\begin{abstract}
This paper develops a second family of special functions generated by exact packet refinement. The starting point is the Bernoulli--Hurwitz child system introduced in the first paper of the series. A finite cyclotomic transform is applied to the children of one parent packet. The Taylor germ of every non-DC detail channel is then identified with a triangular transform of cyclotomic Fourier packets of generalized Stieltjes constants. For a nonprincipal Dirichlet character, the same transform recovers the derivatives $L^{(m)}(1,\chi)$ from generalized Stieltjes constants at rational arguments.

Two scalar packet invariants are introduced. The holomorphic detail variance equals both the sum of the non-DC channel products and the mean pairwise square difference of the children. The primitive packet resultant compresses the complete Galois orbit of primitive cyclotomic details into one entire scalar obstruction. At positive integer order it takes rational values for rational packet parameters.

The arithmetic half of the paper studies the positive squarefree packet energy
\[
E_{a,\beta,\chi}(k)
=
\sum_{n\ge1}
\frac{\mu(n)^2\abs{\chi(n)}^2}{n^{2a}}
\left(1-\frac1{n^\beta}\right)^{2k}.
\]
This sequence is proved to be a Hausdorff moment sequence. All finite differences have alternating sign, every finite Hankel matrix is positive definite, and its generating resolvent is analytic and zero-free on $\C\setminus[1,\infty)$. The same energy is also an exact finite difference of the squarefree zeta ratio $\zeta(s)/\zeta(2s)$ with the Euler factors at the conductor removed.

A complete asymptotic law is proved. After normalization, the active packet coordinate $2k/n^\beta$ converges to a Gamma distribution with shape $(2a-1)/\beta$. For characters modulo one fixed modulus, distinct character packets become asymptotically orthogonal with normalized correlation $O(k^{-1/(2\beta)})$. The principal signed M\"obius sector remains outside this theorem. In the zeta specialization, its subpolynomial coherence is equivalent to the Riemann hypothesis.

All figures were generated by the accompanying Python program. The release contains 20 vector figures, matching PNG files, exact tables, 61 built-in checks, 18 independent tests, and the complete source listing.
\end{abstract}

\chapter*{Status and claim boundary}
\addcontentsline{toc}{chapter}{Status and claim boundary}
The paper proves the Stieltjes-germ formulas, the character transform at $s=1$, the detail-variance identity, the primitive-resultant criterion, complete monotonicity, finite Hankel positivity, the zero-free resolvent theorem, the squarefree zeta-ratio formula, the packet-energy asymptotic, the Gamma scaling limit, and the character-decoherence theorem under the hypotheses stated below.

Classical inputs include the Laurent expansion of the Hurwitz zeta function, Dirichlet-character orthogonality, Bernoulli multiplication, the squarefree counting law, partial summation, the method of moments, and the B\'aez-Duarte criterion. The logarithmic quantum Hamiltonian used in the final part was realized in the 2026 quantum-dynamics work of Wei and collaborators~\cite{Wei2026}.

The phrase ``generated special function'' refers to the exact packet construction and its proved laws. It does not by itself establish historical priority for every scalar specialization. The principal M\"obius coherence statement remains open and is equivalent to RH after normalization by the proved diagonal energy law.

\tableofcontents
\listoffigures

\InlinePart{Cyclotomic germs of Bernoulli refinement}

\FlowChapter{Bernoulli--Hurwitz children}

The preceding paper introduced the entire Bernoulli interpolation
\[
\Bfrak(z;a)=-z\zeta(1-z,a)
\]
and the exact packet family
\[
\Bfrak_{M,r}(z;a)
=
-zM^{z-1}\zeta\!\left(1-z,\frac{a+r}{M}\right).
\]
The reconstruction and refinement laws are
\[
\Bfrak(z;a)=\sum_{r=0}^{M-1}\Bfrak_{M,r}(z;a)
\]
and
\[
\Bfrak_{M,r}(z;a)
=
\sum_{\ell=0}^{L-1}
\Bfrak_{LM,r+\ell M}(z;a).
\]
These formulas were proved in \cite{KomorebiSF1}. The present paper asks what local and scalar information is created by one refinement step.

Fix integers $M\ge1$, $L\ge2$, and $0\le r<M$. Set
\[
Q=LM,
\qquad
q_\ell=\frac{a+r+\ell M}{LM},
\qquad 0\le\ell<L.
\]
The children of the parent packet are
\[
b_\ell(z;a)
=
\Bfrak_{Q,r+\ell M}(z;a)
=
-zQ^{z-1}\zeta(1-z,q_\ell).
\]
Let
\[
\omega_L=\e^{2\pi\ii/L}.
\]
The unitary child transform is
\[
\boxed{
W_j(z;a)
=
\frac1{\sqrt L}
\sum_{\ell=0}^{L-1}
\omega_L^{-j\ell}b_\ell(z;a),
\qquad 0\le j<L.
}
\]
The DC channel is the parent packet divided by $\sqrt L$:
\[
W_0(z;a)=\frac1{\sqrt L}\Bfrak_{M,r}(z;a).
\]
The remaining channels record the exact detail created by the refinement.

\FigureWide{01_theory_pipeline}{The second special-function construction begins with one Bernoulli parent and its exact children.}{The finite child transform produces local Stieltjes packets, scalar detail invariants, and a positive arithmetic energy family.}

\FlowChapter{Cyclotomic Stieltjes packet theorem}

The generalized Stieltjes constants are defined by the Laurent expansion
\[
\zeta(s,q)
=
\frac1{s-1}
+
\sum_{m=0}^{\infty}
\frac{(-1)^m}{m!}\gamma_m(q)(s-1)^m.
\]
Substituting $s=1-z$ yields
\[
\zeta(1-z,q)
=
-\frac1z
+
\sum_{m=0}^{\infty}\gamma_m(q)\frac{z^m}{m!}.
\]

\begin{definition}[Cyclotomic Stieltjes packet]
For $0\le j<L$ and $m\ge0$, define
\[
\boxed{
G_{j,m}(a)
=
\frac1{\sqrt L}
\sum_{\ell=0}^{L-1}
\omega_L^{-j\ell}\gamma_m(q_\ell).
}
\]
\end{definition}

\begin{theorem}[Complete local packet expansion]
For every channel $j$,
\[
\boxed{
W_j(z;a)
=
\frac{\e^{z\log Q}}{Q}
\left[
\sqrt L\,\delta_{j0}
-
\sum_{m=0}^{\infty}
G_{j,m}(a)\frac{z^{m+1}}{m!}
\right].
}
\]
In particular, every non-DC detail has a zero at the removable Bernoulli point:
\[
\boxed{W_j(0;a)=0\qquad (j\ne0).}
\]
For $j\ne0$ and $n\ge1$,
\[
\boxed{
W_j^{(n)}(0;a)
=
-\frac{n}{Q}
\sum_{m=0}^{n-1}
\binom{n-1}{m}
(\log Q)^{n-1-m}G_{j,m}(a).
}
\]
\end{theorem}

\begin{proof}
Insert the Laurent expansion into the child definition:
\[
\begin{aligned}
b_\ell(z;a)
&=-zQ^{z-1}
\left[
-\frac1z+
\sum_{m\ge0}\gamma_m(q_\ell)\frac{z^m}{m!}
\right]\\
&=\frac{\e^{z\log Q}}{Q}
\left[
1-
\sum_{m\ge0}\gamma_m(q_\ell)\frac{z^{m+1}}{m!}
\right].
\end{aligned}
\]
Apply the finite Fourier transform over $\ell$. The transform of the constant vector is $\sqrt L$ in the DC channel and zero in every detail channel. This proves the first formula and the zero at $z=0$.

For a non-DC channel, multiply the two power series
\[
\e^{z\log Q}
\quad\hbox{and}\quad
-\sum_{m\ge0}G_{j,m}(a)\frac{z^{m+1}}{m!}.
\]
The coefficient of $z^n$ is
\[
-\frac1Q
\sum_{m=0}^{n-1}
\frac{G_{j,m}(a)}{m!}
\frac{(\log Q)^{n-1-m}}{(n-1-m)!}.
\]
Multiplication by $n!$ gives the derivative formula.
\end{proof}

\FigureWide{02_stieltjes_packet_matrix}{The real and imaginary parts of the first cyclotomic Stieltjes packets for a four-channel refinement.}{The channels contain distinct local Taylor data even though they reconstruct one Bernoulli parent.}

\begin{example}[The first Bernoulli--Euler detail coefficient]
Take $L=2$, $M=1$, $r=0$, and $a=1$. Then
\[
q_0=\frac12,
\qquad
q_1=1.
\]
The first non-DC packet is
\[
G_{1,0}(1)
=
\frac{\gamma_0(1/2)-\gamma_0(1)}{\sqrt2}.
\]
Since $\gamma_0(q)=-\psi(q)$ and
\[
\psi(1/2)=-\gamma-2\log2,
\qquad
\psi(1)=-\gamma,
\]
we have
\[
G_{1,0}(1)=\sqrt2\log2.
\]
The derivative theorem gives
\[
\boxed{
W_1'(0;1)
=-\frac{\log2}{\sqrt2}.
}
\]
The first four derivatives computed from the packet formula and from direct high-precision differentiation are shown below.
\end{example}

\begin{center}
\begin{tabular}{rrrr}
\toprule
$n$ & packet formula & direct derivative & absolute error\\
\midrule
1 & -0.490129071734 & -0.490129071734 & 2.62e-62\\
2 & 0.226088771874 & 0.226088771874 & 5.55e-17\\
3 & 0.138676210506 & 0.138676210506 & 2.27e-62\\
4 & 0.026626672189 & 0.026626672189 & 1.11e-16\\
\bottomrule
\end{tabular}
\end{center}

\FigureWide{03_two_channel_germ}{The first two-channel detail and its tangent at the removable point.}{The linear coefficient is the exact constant $-(\log2)/\sqrt2$.}

\FlowChapter{Character packets and derivatives at one}

The additive child transform has a multiplicative counterpart. Let $\chi$ be a nonprincipal Dirichlet character modulo $q$. The classical Hurwitz representation is
\[
q^sL(s,\chi)
=
\sum_{r=1}^{q}
\chi(r)\zeta\!\left(s,\frac rq\right).
\]
The pole terms cancel because
\[
\sum_{r=1}^{q}\chi(r)=0.
\]

\begin{definition}[Character Stieltjes packet]
Define
\[
S_m(\chi)
=
\sum_{r=1}^{q}
\chi(r)\gamma_m\!\left(\frac rq\right).
\]
\end{definition}

\begin{theorem}[Character transform at $s=1$]
For every $m\ge0$,
\[
\boxed{
S_m(\chi)
=
(-1)^m q
\sum_{h=0}^{m}
\binom{m}{h}
(\log q)^{m-h}
L^{(h)}(1,\chi).
}
\]
\end{theorem}

\begin{proof}
Write $s=1+w$. The pole-free left side has expansion
\[
\sum_{m\ge0}
\frac{(-1)^m}{m!}S_m(\chi)w^m.
\]
The right side is
\[
q\e^{w\log q}
\sum_{h\ge0}
\frac{L^{(h)}(1,\chi)}{h!}w^h.
\]
Equate the coefficient of $w^m$ and multiply by $m!$.
\end{proof}

The first two cases are
\[
\boxed{
\sum_{r=1}^{q}
\chi(r)\gamma_0\!\left(\frac rq\right)
=qL(1,\chi)
}
\]
and
\[
\boxed{
\sum_{r=1}^{q}
\chi(r)\gamma_1\!\left(\frac rq\right)
=-q\left[L'(1,\chi)+(\log q)L(1,\chi)\right].
}
\]
The packet family therefore joins the special values $L(1-n,\chi)$, carried by the integer Bernoulli samples, to the local derivatives $L^{(m)}(1,\chi)$, carried by the Stieltjes germ.

\begin{center}
\begin{tabular}{rrrr}
\toprule
$m$ & Stieltjes packet sum & $L$-derivative formula & error\\
\midrule
0 & 3.1415926536 & 3.1415926536 & 0\\
1 & -5.1267774478 & -5.1267774482 & 4.31e-10\\
2 & 7.5603277998 & 7.5603278010 & 1.19e-09\\
3 & -10.6337651115 & -10.6337651140 & 2.48e-09\\
\bottomrule
\end{tabular}
\end{center}

\FigureWide{04_character_L_derivatives}{The nonprincipal character modulo four gives the same values from generalized Stieltjes packets and from derivatives of the Dirichlet beta function at one.}{The visible numerical differences come from independent differentiation of the $L$-function implementation.}

\InlinePart{Scalar invariants of a complete refinement}

\FlowChapter{Holomorphic detail variance}

Let $b_0,\ldots,b_{L-1}$ be the child packets at fixed $(z,a)$ and let $W_j$ be their normalized finite Fourier transform.

\begin{definition}[Holomorphic detail variance]
Define
\[
\boxed{
\Vcal_L
=
\sum_{\ell=0}^{L-1}b_\ell^2
-
\frac1L
\left(\sum_{\ell=0}^{L-1}b_\ell\right)^2.
}
\]
\end{definition}

\begin{theorem}[Three exact forms]
The detail variance satisfies
\[
\boxed{
\Vcal_L
=
\sum_{j=1}^{L-1}W_jW_{L-j}
}
\]
and
\[
\boxed{
\Vcal_L
=
\frac1L
\sum_{0\le\ell<m<L}
(b_\ell-b_m)^2.
}
\]
For real $z$ and real $a$, the children are real and
\[
\boxed{
\Vcal_L
=
\sum_{j=1}^{L-1}\abs{W_j}^2\ge0.
}
\]
It vanishes precisely when every child value is equal.
\end{theorem}

\begin{proof}
Fourier orthogonality gives
\[
\sum_{j=0}^{L-1}W_jW_{L-j}
=
\sum_{\ell=0}^{L-1}b_\ell^2.
\]
The DC term is
\[
W_0^2
=
\frac1L\left(\sum_\ell b_\ell\right)^2.
\]
Subtracting it proves the first identity. Expanding all pairwise squares gives
\[
\sum_{\ell<m}(b_\ell-b_m)^2
=
L\sum_\ell b_\ell^2
-
\left(\sum_\ell b_\ell\right)^2,
\]
which proves the second identity. For real children, $W_{L-j}=\overline{W_j}$, so the first identity becomes a sum of squares.
\end{proof}

\FigureWide{05_detail_variance_identity}{The same refinement viewed in child coordinates and cyclotomic coordinates.}{The non-DC transform channels contain exactly the scalar variance lost by replacing all children by their parent mean.}

\FlowChapter{Bernoulli--Euler and three-channel examples}

\begin{proposition}[Euler detail as variance square root]
In the two-channel refinement of the Bernoulli parent,
\[
W_1(n;a)
=-\frac{n}{2\sqrt2}E_{n-1}(a).
\]
Consequently,
\[
\boxed{
\Vcal_2(n;a)
=
\frac{n^2}{8}E_{n-1}(a)^2.
}
\]
\end{proposition}

\begin{proof}
The two children at integer order are
\[
b_0=2^{n-1}B_n(a/2),
\qquad
b_1=2^{n-1}B_n((a+1)/2).
\]
The standard Bernoulli--Euler relation gives
\[
b_0-b_1=-\frac n2E_{n-1}(a).
\]
Divide by $\sqrt2$ to obtain $W_1$. Since $\Vcal_2=W_1^2$, the variance formula follows.
\end{proof}

\FigureWide{06_bernoulli_euler_variance}{The nonnegative two-channel variance for three even Bernoulli orders.}{The Euler polynomial controls the complete detail energy of the Bernoulli split.}

\begin{example}[A three-channel rational variance]
Take $L=3$, $M=1$, $r=0$, $a=0$, and $n=4$. The exact child values are
\[
(b_0,b_1,b_2)
=
\left(-\frac9{10},\frac{13}{30},\frac{13}{30}\right).
\]
Therefore
\[
\begin{aligned}
\Vcal_3
&=\frac13\left[
\left(-\frac43\right)^2
+0^2
+\left(\frac43\right)^2
\right]\\
&=\boxed{\frac{32}{27}}.
\end{aligned}
\]
The two primitive channels are conjugate and satisfy
\[
W_1W_2=\frac{16}{27}.
\]
\end{example}

\FigureWide{07_three_channel_variance}{The three exact order-four Bernoulli children.}{Two children coincide, while the remaining difference produces the rational detail variance $32/27$.}

\FlowChapter{Primitive packet resultants}

Define the child polynomial
\[
P_{M,r}^{(L)}(X;z,a)
=
\sum_{\ell=0}^{L-1}b_\ell(z;a)X^\ell.
\]
At a root of unity,
\[
P_{M,r}^{(L)}(\omega_L^{-j};z,a)
=\sqrt L\,W_j(z;a).
\]

\begin{definition}[Primitive packet resultant]
Let $\Phi_L$ be the $L$th cyclotomic polynomial. Define
\[
\boxed{
\Ncal_{M,r}^{(L)}(z;a)
=
\Res_X\left(\Phi_L(X),P_{M,r}^{(L)}(X;z,a)\right).
}
\]
Equivalently,
\[
\Ncal_{M,r}^{(L)}(z;a)
=
\prod_{\substack{1\le j\le L\\(j,L)=1}}
P_{M,r}^{(L)}(\omega_L^j;z,a).
\]
\end{definition}

\begin{theorem}[Primitive collapse criterion]
The primitive resultant is entire in $z$. It vanishes exactly when at least one primitive cyclotomic detail channel vanishes. If $a\in\Q$ and $n\in\N$, then
\[
\boxed{
\Ncal_{M,r}^{(L)}(n;a)\in\Q.
}
\]
\end{theorem}

\begin{proof}
The resultant is a polynomial in the coefficients $b_\ell(z;a)$, so it is entire. Since $\Phi_L$ is monic, the product formula for the resultant is the product of $P$ over all primitive $L$th roots. Therefore it is zero exactly when one primitive evaluation is zero. At integer order and rational $a$, every child equals a rational multiple of a Bernoulli-polynomial value at a rational point, hence is rational. The resultant of two rational polynomials is rational.
\end{proof}

\FigureWide{08_primitive_resultants}{Primitive packet resultants at integer orders for three cyclotomic depths.}{The plotted scalar is a Galois-invariant compression of all primitive detail channels.}

\InlinePart{Positive arithmetic packet functions}

\FlowChapter{The squarefree packet energy}

Let $a>1/2$, $\beta>0$, and let $\chi$ be a Dirichlet character modulo $q$. Define
\[
\boxed{
E_{a,\beta,\chi}(k)
=
\sum_{n=1}^{\infty}
\frac{\mu(n)^2\abs{\chi(n)}^2}{n^{2a}}
\left(1-\frac1{n^\beta}\right)^{2k}.
}
\]
Set
\[
w_n=\frac{\mu(n)^2\abs{\chi(n)}^2}{n^{2a}},
\qquad
x_n=\left(1-\frac1{n^\beta}\right)^2.
\]
Then
\[
E_{a,\beta,\chi}(k)=\sum_nw_nx_n^k.
\]
The support lies in $[0,1)$ and the total mass is finite because $a>1/2$.

\begin{theorem}[Complete monotonicity]
For all integers $m,k\ge0$,
\[
\boxed{
(-1)^m\Delta^mE_{a,\beta,\chi}(k)
=
\sum_nw_nx_n^k(1-x_n)^m
\ge0.
}
\]
Thus $E_{a,\beta,\chi}$ is a Hausdorff moment sequence. In particular it is decreasing and log-convex:
\[
E(k)^2\le E(k-1)E(k+1).
\]
\end{theorem}

\begin{proof}
Since $\Delta x^k=x^k(x-1)$, repeated finite differences give
\[
\Delta^m x^k=x^k(x-1)^m.
\]
Sum against the nonnegative weights $w_n$. Log-convexity is the Cauchy--Schwarz inequality applied to $x_n^{(k-1)/2}$ and $x_n^{(k+1)/2}$ in the weighted sequence space.
\end{proof}

\FigureWide{09_complete_monotonicity}{The first nine alternating finite-difference layers of the principal packet energy.}{Every displayed quantity is nonnegative, as required by the exact Hausdorff moment representation.}

\FlowChapter{Hankel positivity and the zero-free resolvent}

\begin{theorem}[Strict finite Hankel positivity]
For every $d\ge0$, the Hankel matrix
\[
H_d=\left(E(i+j)\right)_{0\le i,j\le d}
\]
is positive definite.
\end{theorem}

\begin{proof}
For a nonzero vector $c=(c_0,\ldots,c_d)$,
\[
\begin{aligned}
c^*H_dc
&=\sum_{i,j}c_i\overline{c_j}E(i+j)\\
&=\sum_nw_n\left|\sum_{i=0}^{d}c_ix_n^i\right|^2.
\end{aligned}
\]
The support contains infinitely many distinct $x_n$ with positive weight. A nonzero polynomial of degree at most $d$ cannot vanish at all of them. Hence the sum is strictly positive.
\end{proof}

The smallest computed eigenvalues eventually reach floating-point roundoff. The theorem is exact and does not depend on those numerical signs.

\begin{center}
\begin{tabular}{rrr}
\toprule
size & smallest eigenvalue & largest eigenvalue\\
\midrule
2 & 6.345e-04 & 1.075e-01\\
3 & 1.865e-06 & 1.197e-01\\
4 & 6.036e-10 & 1.253e-01\\
5 & 8.948e-14 & 1.282e-01\\
6 & -2.410e-17 & 1.297e-01\\
7 & -2.218e-17 & 1.307e-01\\
8 & -3.034e-17 & 1.313e-01\\
9 & -3.447e-17 & 1.317e-01\\
10 & -3.586e-17 & 1.320e-01\\
\bottomrule
\end{tabular}
\end{center}

\FigureWide{10_hankel_positivity}{The largest and smallest eigenvalues of finite packet-energy Hankel matrices.}{The tiny negative values at the largest sizes are at floating-point roundoff. The exact Gram representation is strictly positive.}

\begin{definition}[Packet-energy resolvent]
Define
\[
\boxed{
\Ecal_{a,\beta,\chi}(z)
=
\sum_{k=0}^{\infty}E_{a,\beta,\chi}(k)z^k.
}
\]
For $\abs z<1$,
\[
\Ecal_{a,\beta,\chi}(z)
=
\sum_n\frac{w_n}{1-zx_n}.
\]
\end{definition}

\begin{theorem}[Zero-free slit plane]
The resolvent extends analytically to
\[
\C\setminus[1,\infty)
\]
and satisfies
\[
\boxed{
\Ecal_{a,\beta,\chi}(z)\ne0
\qquad
(z\in\C\setminus[1,\infty)).
}
\]
\end{theorem}

\begin{proof}
Uniform convergence on compact subsets of the slit plane follows from $\sum_nw_n<\infty$ and a uniform lower bound on $\abs{1-zx}$ for $x\in[0,1]$.

For $z\ne0$, set $u=1/z$ and
\[
S(u)=\sum_n\frac{w_n}{u-x_n}.
\]
Then
\[
\Ecal(z)=uS(u).
\]
If $\Im u>0$, then
\[
\Im S(u)
=-\Im u\sum_n\frac{w_n}{\abs{u-x_n}^2}<0.
\]
The sign reverses in the lower half-plane. If $u>1$, then every summand is positive. If $u<0$, every summand is negative. These cases cover the image of the slit domain under $u=1/z$. Therefore $S(u)$ and $\Ecal(z)$ never vanish. At $z=0$, $\Ecal(0)=\sum_nw_n>0$.
\end{proof}

\FigureWide{11_zero_free_resolvent}{The magnitude of the principal packet-energy resolvent on the slit plane.}{The excluded positive-real cut contains the poles and their accumulation at one. The theorem rules out zeros everywhere else.}

\FlowChapter{A finite-difference zeta-ratio coordinate system}

Because $\abs{\chi(n)}^2$ is the indicator of $(n,q)=1$, define
\[
D_q(s)
=
\sum_{\substack{n\ge1\\(n,q)=1}}
\frac{\mu(n)^2}{n^s}.
\]
Its Euler product is
\[
\boxed{
D_q(s)
=
\frac{\zeta(s)}{\zeta(2s)}
\prod_{p\mid q}
\left(1+p^{-s}\right)^{-1}.
}
\]

\begin{theorem}[Exact zeta-ratio finite difference]
For every integer $k\ge0$,
\[
\boxed{
E_{a,\beta,\chi}(k)
=
\sum_{j=0}^{2k}
(-1)^j\binom{2k}{j}
D_q(2a+\beta j).
}
\]
\end{theorem}

\begin{proof}
Expand the filter by the binomial theorem:
\[
\left(1-n^{-\beta}\right)^{2k}
=
\sum_{j=0}^{2k}
(-1)^j\binom{2k}{j}n^{-\beta j}.
\]
Absolute convergence permits interchange with the squarefree Dirichlet series. The inner sum is $D_q(2a+\beta j)$.
\end{proof}

\begin{center}
\begin{tabular}{rrrr}
\toprule
$k$ & direct energy & zeta-ratio difference & truncation error\\
\midrule
1 & 4.7804610e-02 & 4.7804610e-02 & 7.494e-16\\
2 & 3.0200677e-02 & 3.0200677e-02 & 4.406e-16\\
3 & 1.9767515e-02 & 1.9767515e-02 & 3.227e-16\\
4 & 1.3470094e-02 & 1.3470094e-02 & 2.481e-16\\
5 & 9.5803728e-03 & 9.5803728e-03 & 1.665e-16\\
6 & 7.1098189e-03 & 7.1098189e-03 & 1.639e-16\\
7 & 5.4893011e-03 & 5.4893011e-03 & 1.266e-16\\
8 & 4.3883096e-03 & 4.3883096e-03 & 8.153e-17\\
\bottomrule
\end{tabular}
\end{center}

\FigureWide{12_zeta_ratio_coordinates}{Direct squarefree summation and the exact finite difference of $\zeta(s)/\zeta(2s)$ agree.}{The small displayed difference is only the truncation error of the direct numerical sum.}

\InlinePart{Asymptotic packet geometry}

\FlowChapter{The diagonal energy law}

Let
\[
A_q(x)
=
\sum_{\substack{n\le x\\(n,q)=1}}\mu(n)^2.
\]
The elementary squarefree counting estimate is
\[
A_q(x)=\delta_qx+O_q(\sqrt x),
\]
where
\[
\boxed{
\delta_q
=
\frac1{\zeta(2)}
\prod_{p\mid q}\frac{p}{p+1}.
}
\]
Set
\[
\alpha=\frac{2a-1}{\beta}.
\]

\begin{theorem}[Packet-energy asymptotic]
For $a>1/2$ and $\beta>0$,
\[
\boxed{
E_{a,\beta,\chi}(k)
\sim
\frac{\delta_q}{\beta}
\Gamma(\alpha)
(2k)^{-\alpha}.
}
\]
More precisely, the error produced by the $O_q(\sqrt x)$ remainder is
\[
O_{a,\beta,q}\!\left(k^{-(2a-1/2)/\beta}\right).
\]
\end{theorem}

\begin{proof}
Use Stieltjes summation with
\[
f_k(x)=x^{-2a}(1-x^{-\beta})^{2k}.
\]
The density term contributes
\[
\delta_q\int_1^\infty f_k(x)\dd x.
\]
With $y=x^{-\beta}$,
\[
\begin{aligned}
\int_1^\infty f_k(x)\dd x
&=\frac1\beta
\int_0^1
 y^{\alpha-1}(1-y)^{2k}\dd y\\
&=\frac1\beta
\mathrm B(\alpha,2k+1).
\end{aligned}
\]
The beta-function asymptotic gives
\[
\mathrm B(\alpha,2k+1)
\sim
\Gamma(\alpha)(2k)^{-\alpha}.
\]

Write the squarefree counting remainder as $R_q(x)=O_q(\sqrt x)$. Partial summation bounds its contribution by a constant times
\[
\int_1^\infty\sqrt x\abs{f_k'(x)}\dd x.
\]
Scaling $x=k^{1/\beta}y$ shows that this integral is
\[
O\!\left(k^{-(2a-1/2)/\beta}\right),
\]
which is smaller than the main term by $k^{-1/(2\beta)}$.
\end{proof}

For $a=\beta=2$ and the principal character,
\[
\boxed{
E_k
\sim
\frac{3}{2^{5/2}\pi^{3/2}}k^{-3/2}.
}
\]

\begin{center}
\begin{tabular}{rrrr}
\toprule
$k$ & $E_k$ & $k^{3/2}E_k$ & relative error\\
\midrule
100 & 9.397187e-05 & 0.0939718675 & 1.332e-02\\
300 & 1.844163e-05 & 0.0958255236 & 6.143e-03\\
1000 & 3.028409e-06 & 0.0957666957 & 5.525e-03\\
3000 & 5.760934e-07 & 0.0946618062 & 6.076e-03\\
10000 & 9.539419e-08 & 0.0953941855 & 1.614e-03\\
30000 & 1.833371e-08 & 0.0952647535 & 2.551e-04\\
100000 & 3.007053e-09 & 0.0950913727 & 1.565e-03\\
\bottomrule
\end{tabular}
\end{center}

\FigureWide{13_energy_asymptotic}{The scaled principal energy approaches the exact constant $3/(2^{5/2}\pi^{3/2})$.}{The final relative error in the finite computation is about $0.16$ percent.}

\FlowChapter{Universal Gamma scaling}

Normalize the packet weights into a probability distribution:
\[
\mathbb P_k(n)
=
\frac{
\mu(n)^2\abs{\chi(n)}^2n^{-2a}(1-n^{-\beta})^{2k}
}{E_{a,\beta,\chi}(k)}.
\]
Define
\[
\boxed{
U_k=\frac{2k}{n^\beta}.
}
\]

\begin{theorem}[Gamma scaling limit]
As $k\to\infty$,
\[
\boxed{
U_k\Longrightarrow\operatorname{Gamma}(\alpha,1),
\qquad
\alpha=\frac{2a-1}{\beta}.
}
\]
Equivalently, for every bounded continuous $\varphi$,
\[
\lim_{k\to\infty}
\sum_n\mathbb P_k(n)
\varphi\!\left(\frac{2k}{n^\beta}\right)
=
\frac1{\Gamma(\alpha)}
\int_0^\infty
\varphi(u)u^{\alpha-1}\e^{-u}\dd u.
\]
\end{theorem}

\begin{proof}
For every integer $m\ge0$,
\[
\E_k[U_k^m]
=(2k)^m
\frac{E_{a+\beta m/2,\beta,\chi}(k)}{E_{a,\beta,\chi}(k)}.
\]
Apply the packet-energy asymptotic to numerator and denominator. The squarefree density cancels and gives
\[
\lim_{k\to\infty}\E_k[U_k^m]
=
\frac{\Gamma(\alpha+m)}{\Gamma(\alpha)}.
\]
These are the moments of $\operatorname{Gamma}(\alpha,1)$. The Gamma distribution is moment-determinate. The first moments remain bounded, so the family is tight. The method of moments gives weak convergence.
\end{proof}

For the principal RH parameters $a=\beta=2$,
\[
\alpha=\frac32.
\]
The active packet window therefore has a complete limiting profile, not only the rough location $n\asymp\sqrt k$.

\begin{center}
\resizebox{\textwidth}{!}{\begin{tabular}{rrrrrrr}
\toprule
$k$ & mean & variance & $q_{0.1}$ & $q_{0.5}$ & $q_{0.9}$ & KS\\
\midrule
300 & 1.54561 & 1.54638 & 0.31798 & 1.18207 & 3.01649 & 0.05780\\
3000 & 1.50231 & 1.59381 & 0.29310 & 1.17179 & 3.29253 & 0.04726\\
30000 & 1.49275 & 1.48729 & 0.29110 & 1.20548 & 3.08603 & 0.02256\\
100000 & 1.50160 & 1.52410 & 0.29473 & 1.16989 & 3.10136 & 0.01560\\
\bottomrule
\end{tabular}}
\end{center}

\FigureWide{14_gamma_scaling}{The empirical active-packet distributions approach the Gamma$(3/2,1)$ density.}{The horizontal coordinate is the intrinsic inverse scale $U_k=2k/n^2$.}

\FigureWide{15_gamma_convergence}{The weighted distribution approaches its Gamma limit as the packet order increases.}{The plotted distance is a finite numerical certificate rather than the proof, which follows from the moment asymptotics.}

\FlowChapter{Character packet decoherence}

Let $\chi$ and $\psi$ be Dirichlet characters modulo the same integer $q$. Define the packet Gram entry
\[
\boxed{
G_{\chi,\psi}(k)
=
\sum_{n=1}^{\infty}
\frac{\mu(n)^2\chi(n)\overline{\psi(n)}}{n^{2a}}
\left(1-\frac1{n^\beta}\right)^{2k}.
}
\]
The diagonal is the positive packet energy.

\begin{lemma}[Squarefree character cancellation]
If $\xi$ is a nonprincipal Dirichlet character modulo $q$, then
\[
\boxed{
\sum_{n\le X}\mu(n)^2\xi(n)=O_q(\sqrt X).
}
\]
\end{lemma}

\begin{proof}
Use
\[
\mu(n)^2=\sum_{d^2\mid n}\mu(d).
\]
Then
\[
\sum_{n\le X}\mu(n)^2\xi(n)
=
\sum_{d\le\sqrt X}
\mu(d)\xi(d^2)
\sum_{m\le X/d^2}\xi(m).
\]
A nonprincipal character has zero sum over a complete period, so every partial sum is $O(q)$. There are $O(\sqrt X)$ outer terms.
\end{proof}

\begin{theorem}[Character decoherence]
If $\chi\ne\psi$, then
\[
\boxed{
G_{\chi,\psi}(k)
=
O_{a,\beta,q}
\left(k^{-(2a-1/2)/\beta}\right).
}
\]
Consequently,
\[
\boxed{
\frac{G_{\chi,\psi}(k)}
{\sqrt{G_{\chi,\chi}(k)G_{\psi,\psi}(k)}}
=
O_{a,\beta,q}\left(k^{-1/(2\beta)}\right).
}
\]
\end{theorem}

\begin{proof}
The product $\xi=\chi\overline\psi$ is nonprincipal. Let
\[
A_\xi(X)=\sum_{n\le X}\mu(n)^2\xi(n)=O_q(\sqrt X).
\]
Apply partial summation to
\[
f_k(x)=x^{-2a}(1-x^{-\beta})^{2k}.
\]
The resulting bound is a constant times
\[
\int_1^\infty\sqrt x\abs{f_k'(x)}\dd x
=O\!\left(k^{-(2a-1/2)/\beta}\right).
\]
The two diagonal entries are asymptotic to positive constants times $k^{-(2a-1)/\beta}$. Dividing gives the normalized rate.
\end{proof}

For $a=\beta=2$, the normalized decay rate is
\[
\boxed{k^{-1/4}.}
\]

\FigureWide{16_character_gram}{The absolute normalized Gram matrix for the four characters modulo five at two packet orders.}{The diagonal remains one while the off-diagonal channels decay.}

\begin{center}
\begin{tabular}{rr}
\toprule
$k$ & largest normalized off-diagonal entry\\
\midrule
100 & 0.09807746\\
300 & 0.07180785\\
1000 & 0.05948305\\
3000 & 0.05265039\\
10000 & 0.01447532\\
30000 & 0.01243919\\
100000 & 0.01185877\\
\bottomrule
\end{tabular}
\end{center}

\FigureWide{17_character_decoherence}{The largest off-diagonal correlation in the character packet bank.}{The theorem gives the general $k^{-1/(2\beta)}$ upper rate. The finite values need not be monotone at every intermediate order.}

\InlinePart{Quantum realization and the remaining principal obstruction}

\FlowChapter{A logarithmic packet Hamiltonian}

Work in
\[
\mathcal H=\ell^2(\{2,3,\ldots\}).
\]
For $a>1$, define the amplitude vector
\[
\ket{\phi_a}
=
\sum_{n\ge2}n^{-a/2}\ket n.
\]
For $a>1/2$, define the diagonal-power vector
\[
\ket{\psi_a}
=
\sum_{n\ge2}n^{-a}\ket n.
\]
Let
\[
H\ket n=\log n\,\ket n,
\]
\[
K_\beta\ket n
=-\log(1-n^{-\beta})\ket n,
\]
and
\[
\Mcal_\chi\ket n
=\mu(n)\chi(n)\ket n.
\]
The three operators are diagonal and commute.

\begin{proposition}[Quantum transfer amplitude]
The generalized signed transfer is
\[
\boxed{
\Ccal_{a,\beta,\chi}(k,t)
=
\bra{\phi_a}
\Mcal_\chi
\e^{-kK_\beta}
\e^{-\ii tH}
\ket{\phi_a}.
}
\]
Its positive diagonal power is
\[
\boxed{
E_{a,\beta,\chi}(k)
=
\bra{\psi_a}
\Mcal_\chi^*\Mcal_\chi
\e^{-2kK_\beta}
\ket{\psi_a}.
}
\]
It is also the long-time mean power of the signed transfer:
\[
\boxed{
E_{a,\beta,\chi}(k)
=
\lim_{T\to\infty}
\frac1{2T}
\int_{-T}^{T}
\abs{\Ccal_{a,\beta,\chi}(k,t)}^2\dd t.
}
\]
\end{proposition}

\begin{proof}
The first matrix element has number-basis factors
\[
n^{-a},
\qquad
\mu(n)\chi(n),
\qquad
(1-n^{-\beta})^k,
\qquad
n^{-\ii t},
\]
which give the signed transfer. The second matrix element has factors
\[
n^{-2a},
\qquad
\mu(n)^2\abs{\chi(n)}^2,
\qquad
(1-n^{-\beta})^{2k},
\]
which give the positive packet energy. Finally, expand the squared modulus of the transfer. The time average of
\(
\e^{-\ii t(\log n-\log m)}
\)
vanishes when $n\ne m$ and equals one when $n=m$. Only the diagonal terms remain.
\end{proof}

The logarithmic Hamiltonian is not introduced only as formal notation. Wei and collaborators construct a quantum system with spectrum $E_n=\log n$, thermal weights $n^{-\beta}$, and dynamical phases $n^{-\ii t}$~\cite{Wei2026}. Their observable directly encodes a zeta amplitude. The present packet system uses the same logarithmic spectral coordinate, then adds the squarefree or character sign and the scale filter $\e^{-kK_\beta}$.

\FigureWide{18_quantum_packet_pipeline}{The positive and signed quantum packet channels.}{The logarithmic Hamiltonian supplies real-time phases. The scale Hamiltonian localizes the active number packets.}

\FlowChapter{The principal M\"obius coherence problem}

In the principal zeta specialization, set
\[
a=2,
\qquad
\beta=2,
\qquad
\chi=1.
\]
Then
\[
c_k
=
\sum_{n=1}^{\infty}
\frac{\mu(n)}{n^2}
\left(1-\frac1{n^2}\right)^k
\]
which is the B\'aez-Duarte coefficient. Define
\[
\Gamma_k
=
\frac{\abs{c_k}^2}{E_k}.
\]
The diagonal theorem gives
\[
E_k\asymp k^{-3/2}.
\]
The B\'aez-Duarte criterion gives
\[
\mathrm{RH}
\Longleftrightarrow
c_k=O_\varepsilon(k^{-3/4+\varepsilon})
\quad\hbox{for every }\varepsilon>0.
\]
Therefore:

\begin{theorem}[Normalized principal coherence criterion]
\[
\boxed{
\mathrm{RH}
\Longleftrightarrow
\Gamma_k=O_\varepsilon(k^\varepsilon)
\quad\hbox{for every }\varepsilon>0.
}
\]
\end{theorem}

\begin{proof}
If RH holds, square the B\'aez-Duarte estimate and divide by $E_k\asymp k^{-3/2}$. This gives $\Gamma_k=O_\varepsilon(k^{2\varepsilon})$, which is the stated form after renaming the arbitrary exponent.

Conversely, if $\Gamma_k=O_\varepsilon(k^\varepsilon)$, then
\[
\abs{c_k}^2
\ll_\varepsilon
k^{-3/2+\varepsilon}.
\]
Taking square roots gives the B\'aez-Duarte estimate after another renaming of the arbitrary exponent.
\end{proof}

The character-decoherence theorem controls interference between distinct character sectors. It does not control the internal M\"obius coherence of the principal sector. That remaining signed estimate carries the RH-scale difficulty.

\begin{conjecture}[Principal packet coherence]
For every $\varepsilon>0$,
\[
\boxed{
\left|
\sum_{n=1}^{\infty}
\frac{\mu(n)}{n^2}
\left(1-\frac1{n^2}\right)^k
\right|^2
\le
C_\varepsilon k^\varepsilon
\sum_{n=1}^{\infty}
\frac{\mu(n)^2}{n^4}
\left(1-\frac1{n^2}\right)^{2k}.
}
\]
\end{conjecture}

This conjecture is equivalent to RH. It is not promoted as a proof.

\FigureWide{19_status_map}{The exact boundary between the proved positive theory and the open signed theory.}{The positive energy and cross-character layers are controlled. The principal M\"obius coherence remains open.}

\FlowChapter{Research questions generated by the calculus}

The results above create several concrete questions that are separate from RH.

\begin{enumerate}[label=\arabic*.]
\item Determine the zero sets and growth order of the primitive packet resultants $\Ncal_{M,r}^{(L)}(z;a)$.
\item Classify the real parameters for which the holomorphic detail variance has only real or nonnegative integer samples.
\item Determine whether the Stieltjes packet matrices $G_{j,m}(a)$ obey finite-rank, total-positivity, or variation-diminishing laws for selected cyclotomic channels.
\item Obtain higher-order asymptotic expansions for $E_{a,\beta,\chi}(k)$ and the Gamma scaling law.
\item Sharpen the elementary $O(k^{-1/(2\beta)})$ character-decoherence rate using stronger character-sum estimates.
\item Construct orthogonal polynomials and continued fractions associated with the positive moment sequence $E(k)$.
\item Determine the spectral measure and boundary behavior of the zero-free resolvent $\Ecal(z)$.
\item Extend the packet Gram theory from Dirichlet characters to automorphic coefficient systems.
\item Decide whether the principal M\"obius coherence can be controlled by a local packet inequality rather than a global zero argument.
\end{enumerate}

\FigureWide{20_function_family_map}{The analytic and arithmetic branches of the second generated family.}{Cyclotomic germs, positive moments, Gamma scaling, and character Gram matrices now belong to one exact refinement system.}

\appendix

\FlowChapter{Formula registry}

\begin{longtable}{L{0.31\textwidth}L{0.61\textwidth}}
\toprule
Object & Exact formula\\
\midrule
\endfirsthead
\toprule
Object & Exact formula\\
\midrule
\endhead
Cyclotomic Stieltjes packet & $G_{j,m}(a)=L^{-1/2}\sum_\ell\omega_L^{-j\ell}\gamma_m(q_\ell)$\\
Detail germ & $W_j(z;a)=Q^{-1}\e^{z\log Q}[\sqrt L\delta_{j0}-\sum_mG_{j,m}z^{m+1}/m!]$\\
Character transform & $S_m(\chi)=(-1)^mq\sum_{h\le m}\binom mh(\log q)^{m-h}L^{(h)}(1,\chi)$\\
Detail variance & $\Vcal_L=\sum b_\ell^2-L^{-1}(\sum b_\ell)^2$\\
Primitive resultant & $\Ncal=\Res(\Phi_L,P)$\\
Packet energy & $E(k)=\sum_nw_nx_n^k$\\
Complete monotonicity & $(-1)^m\Delta^mE(k)=\sum_nw_nx_n^k(1-x_n)^m$\\
Resolvent & $\Ecal(z)=\sum_nw_n/(1-zx_n)$\\
Zeta-ratio coordinate & $E(k)=\sum_{j=0}^{2k}(-1)^j\binom{2k}{j}D_q(2a+\beta j)$\\
Energy asymptotic & $E(k)\sim(\delta_q/\beta)\Gamma(\alpha)(2k)^{-\alpha}$\\
Gamma coordinate & $U_k=2k/n^\beta\Rightarrow\operatorname{Gamma}(\alpha,1)$\\
Character decoherence & $\Corr(\chi,\psi)=O(k^{-1/(2\beta)})$ for $\chi\ne\psi$\\
\bottomrule
\end{longtable}

\FlowChapter{Numerical certificate tables}

\section*{Packet-energy asymptotic}
\begin{center}
\begin{tabular}{rrrr}
\toprule
$k$ & $E_k$ & $k^{3/2}E_k$ & relative error\\
\midrule
100 & 9.397187e-05 & 0.0939718675 & 1.332e-02\\
300 & 1.844163e-05 & 0.0958255236 & 6.143e-03\\
1000 & 3.028409e-06 & 0.0957666957 & 5.525e-03\\
3000 & 5.760934e-07 & 0.0946618062 & 6.076e-03\\
10000 & 9.539419e-08 & 0.0953941855 & 1.614e-03\\
30000 & 1.833371e-08 & 0.0952647535 & 2.551e-04\\
100000 & 3.007053e-09 & 0.0950913727 & 1.565e-03\\
\bottomrule
\end{tabular}
\end{center}

\section*{Gamma scaling}
\begin{center}
\resizebox{\textwidth}{!}{\begin{tabular}{rrrrrrr}
\toprule
$k$ & mean & variance & $q_{0.1}$ & $q_{0.5}$ & $q_{0.9}$ & KS\\
\midrule
300 & 1.54561 & 1.54638 & 0.31798 & 1.18207 & 3.01649 & 0.05780\\
3000 & 1.50231 & 1.59381 & 0.29310 & 1.17179 & 3.29253 & 0.04726\\
30000 & 1.49275 & 1.48729 & 0.29110 & 1.20548 & 3.08603 & 0.02256\\
100000 & 1.50160 & 1.52410 & 0.29473 & 1.16989 & 3.10136 & 0.01560\\
\bottomrule
\end{tabular}}
\end{center}

\section*{Character decoherence}
\begin{center}
\begin{tabular}{rr}
\toprule
$k$ & largest normalized off-diagonal entry\\
\midrule
100 & 0.09807746\\
300 & 0.07180785\\
1000 & 0.05948305\\
3000 & 0.05265039\\
10000 & 0.01447532\\
30000 & 0.01243919\\
100000 & 0.01185877\\
\bottomrule
\end{tabular}
\end{center}

\section*{Verification count}
\[
\boxed{61/61\text{ checks passed}}
\]
The independent test suite contains 18 additional tests. It rederives the main identities without reading the built-in check ledger.

\FlowChapter{Implementation map}

\begin{longtable}{L{0.40\textwidth}L{0.52\textwidth}}
\toprule
Function & Role\\
\midrule
\endfirsthead
\toprule
Function & Role\\
\midrule
\endhead
\texttt{bernoulli\_packet} & Evaluates one Bernoulli--Hurwitz child\\
\texttt{stieltjes\_packet} & Computes one cyclotomic Stieltjes packet\\
\texttt{detail\_derivative\_from\_stieltjes} & Builds the local Taylor derivatives\\
\texttt{detail\_variance} & Evaluates the scalar non-DC variance\\
\texttt{primitive\_packet\_resultant\_exact} & Computes the exact cyclotomic resultant\\
\texttt{packet\_energy\_values} & Evaluates the positive squarefree energy\\
\texttt{energy\_resolvent} & Evaluates the slit-plane resolvent\\
\texttt{energy\_from\_zeta\_ratio} & Uses the exact finite-difference coordinate\\
\texttt{gamma\_weighted\_sample} & Builds the active inverse-scale distribution\\
\texttt{character\_gram\_matrix} & Constructs the finite character packet Gram matrix\\
\bottomrule
\end{longtable}

\FlowChapter{Complete Python source}

The following program generates every figure, every numerical table, and the complete built-in check ledger used in this paper.

\lstinputlisting{projector_generated_special_functions_ii_listing.py}

\FlowChapter{Reproduction}

From the source directory, run
\begin{lstlisting}[language=bash,numbers=none]
python build_all.py
\end{lstlisting}
The command regenerates the figures and tables, runs the tests, and compiles the paper three times.

\begin{thebibliography}{99}

\bibitem{KomorebiSF1}
Komorebi P.,
\emph{Projector-Generated Special Functions I: Polyphase Differential Lifts, Bernoulli--Lerch Wavelets, Cyclotomic Detail Functions, and M\"obius Coherence}, 2026.

\bibitem{KomorebiBernoulli}
Komorebi P.,
\emph{Bernoulli Transfer-Block Reformulation of Zeta}, proof-development note, 2026.

\bibitem{DLMFBernoulli}
NIST Digital Library of Mathematical Functions,
Chapter 24, Bernoulli and Euler Polynomials.

\bibitem{DLMFZeta}
NIST Digital Library of Mathematical Functions,
Chapter 25, Zeta and Related Functions.

\bibitem{BaezDuarte}
L. B\'aez-Duarte,
A sequential Riesz-like criterion for the Riemann hypothesis,
\emph{International Journal of Mathematics and Mathematical Sciences}, 2005.

\bibitem{HardyWright}
G. H. Hardy and E. M. Wright,
\emph{An Introduction to the Theory of Numbers},
Oxford University Press.

\bibitem{Wei2026}
S. Wei, Y. Zhai, Q. Lu et al.,
The Riemann Hypothesis manifested in dynamical quantum phase transitions,
\emph{Nature Communications}, accepted 18 June 2026,
DOI 10.1038/s41467-026-74935-8.

\end{thebibliography}

\end{document}
